\scanpage{115}
\mspage{2}

On introduit des sommets
\[
(6)\qquad R_r^{i,j}\qquad r\in\mathcal{R},\quad (i,j)\in\operatorname{\unsure{Rep|Coup}}\{0,1,\infty\}\quad\text{i.e. }i,j\in\{0,1,\infty\},\ i\ne j
\]
comme indiqués, \uline{satisfaisant} aux seules relations
\[
(7)\qquad \boxed{R_r^{ij}=R_{\unsure{\sigma_k(r)|\tau_k(r)}}^{i,j}}\quad k\notin\{i,j\}.
\]
On dit que $R_r^{ij}$ est ``proche'' du $R_r^i$ \inserted{sommet de $r$} ($=R_s,R_{\unsure{o|a}},R_f$ suivant que $i=0,1,\infty$), il se trouve sur le côté \deleted{\unsure{du plus proche}} \inserted{qui joint celui-ci} \deleted{\unsure{sommet opposé}} au sommet de type $j$.

On introduit des arcs
\[
(8)\qquad \lambda_r^i:R_r^{i,i+1}\longrightarrow R_r^{i,i+2}\qquad r\in\mathcal{R},\ i\in\{0,1,\infty\}
\]
\[
(9)\qquad a_r^i:R_r^{\deleted{\unsure{i+2}}\unsure{i},i+1}\longrightarrow R_r^{i+1,i+2}\qquad r\in\mathcal{R},\ i\in\{0,1,\infty\}
\]
\margin{\uline{NB} On écrit pour simplifier $i+1$ au lieu du résultat de la permutation circulaire $(0,1,\infty)\longmapsto(1,\infty,0)$, i.e. on identifie $\{0,1,\infty\}$ à $\mathbb{Z}/3\mathbb{Z}\ldots$}
\deleted{qui} \inserted{les $a_r^i$} entourent le disque repéré $r$ en formant un contour fermé dans le sens \deleted{opposé} \inserted{de} l'orientation de $r$ --- \deleted{par contre les} \inserted{\unsure{de = disques}} $\lambda_r^i$ \deleted{\unsure{sont}} \deleted{\unsure{une}} portions d'arc tournant autour des sommets de $r$ de type $i$ dans le sens correspondant à l'orientation de $r$. (Cela \deleted{\unsure{nous}} semble préférable à la convention de notation inverse). On a les seules \deleted{identités} relations
\[
(10)\qquad \boxed{a_r^i=a_{\tau_i(r)}^i}.
\]
\margin{\uline{NB} On désigne par les \deleted{\unsure{ces}} lettres minuscules les flèches obtenues après bouchage de certains trous.}
\margin{(\uline{NB} \unsure{si|i.e.} on bouche les disques $\overline{R}_r^i$, \unsure{où|et} $\lambda_r^i$ devient $=1$)}

\deleted{On désigne}
\uline{NB} On \unsure{suppose} que $X$ \unsure{est} \unsure{condition de trans-} \unsure{réalisée} par sa réalisation conforme canonique, qui est une surface $\mathbb{C}$ analytique réelle donc $C^\infty$, dont on peut considérer son éclaté en les $R_q$ $(q\in S)$ et le découpage de l'éclaté, qui est une surface à bord analytique réelle (mais pas conforme) dont l'intérieur coïncide avec $X\setminus S$, on le note

\mspage{114}
